\scanpage{50}
\mspage{\deleted{\unsure{5|4}}\unsure{4|5}}

\uline{Corollaire} Soit $X$ une catégorie finie, pseudo-cofiltr. \deleted{\unsure{et|est|sép.}}\inserted{(i.e $X^\wedge$ est loc. irréductible)}. Donc $\operatorname{Ord} X$ est un ensemble ordonné cofiltrant, admettant donc un plus petit élément. Soit $U$ l'ouvert de $X$ formé des objets de $X$, qui sont des plus petits objets. \deleted{On suppose} \inserted{Si} $U$ $1$-connexe, \deleted{alors} \deleted{\unsure{$X$|$U$}} est cofiltrant. (Donc si tout objet $0$-connexe de $X$ est $1$-connexe, alors $X$ est cofiltrant.)

En effet, il est immédiat \deleted{que} $U$ étant un ouvert cocofinal dans $X$, \inserted{(et $X$ étant pseudo cofiltrante)} que $X$ est cofiltrant dès que $U$ l'est. D'autre part, si $e\in U$, alors l'inclusion de la sous-catégorie pleine \inserted{$U_0$} de $X$ réduite à $e$ dans $U$ est une équivalence,\textsuperscript{*} ce qui implique que $U_0$ est $1$-connexe, $U$ l'étant, et que $U$ est cofiltrant sss $U_0$ l'est.
% slanted
\margin{Faux ! revoir la \uline{démonstration}}
\deleted{Donc} Il est immédiat aussi, $X$ étant pseudo-cofiltrant, que l'ouvert cocofinal $U$ l'est également, donc $U_0$ (équivalente à $U$) également.

Soit donc $M=\operatorname{End}_X(e)$. On peut appliquer à $M$ la proposition 1, qui implique que $M$ est cofiltrant, i.e. $U_0\simeq B_M$ est cofiltrant, qed

\uline{NB} Au lieu de l'hypothèse $X$ \uline{finie}, il suffisait de supposer que a) $X$ admet un plus petit objet $e$, et b) $\operatorname{End}_X(e)$ est fini.
% archive sheet number: 49
